%% ma: L10-B14-0102
%% lop: 10
%% chuong: 5
%% bai: 14
%% dang: 10.14.8
%% loai: TLN
%% muc: VD
%% dap_an: "1,33"
%% nguon: "chuyen-de-cac-so-dac-trung-cua-mau-so-lieu-khong-ghep-nhom-toan-10-chuong-trinh-moi.pdf (GV. Phan Nhật Linh) (P05), Bài 03 Dạng 2 Phần III Câu 8"
%% kien_thuc: "Lớp 10 Bài 13 (số trung bình), Bài 14 (phương sai, độ lệch chuẩn)"
%% trang_thai: chinh-thuc
%% ghi_chu: "Bổ sung yêu cầu làm tròn đến hàng phần trăm. Tự duyệt theo yêu cầu giáo viên 22/10/2026."
\begin{ex}
Cho bảng điểm kiểm tra cuối học kì 2 môn Hóa học (thang điểm 10) của $40$ học sinh lớp 10A như sau:
\begin{center}
\begin{tabular}{|c|c|c|c|c|c|c|c|c|c|}\hline
Điểm & $3$ & $4$ & $5$ & $6$ & $7$ & $8$ & $9$ & $10$ & Cộng\\\hline
Tần số & $1$ & $1$ & $2$ & $7$ & $12$ & $14$ & $2$ & $1$ & $40$\\\hline
\end{tabular}
\end{center}
Tính độ lệch chuẩn của mẫu số liệu (làm tròn đến hàng phần trăm).
\shortans{1,33}
\loigiai{$\overline x=\dfrac{283}{40}=7{,}075$; $s^2\approx1{,}769$; $s\approx1{,}33$.}
\end{ex}
